\documentclass[a4paper,11pt]{article}
\usepackage[utf8]{inputenc}
\usepackage[T1]{fontenc}
\usepackage[french]{babel}
\usepackage[left=2cm,right=2cm,top=2cm,bottom=2cm]{geometry}
\usepackage{xcolor}
\usepackage{hyperref}
\usepackage{fontawesome5}
\usepackage{parskip}

% --- Couleurs Groupe SII ---
\definecolor{primary}{RGB}{0, 75, 135}
\hypersetup{colorlinks=true, urlcolor=primary}

\begin{document}

% --- EXPÉDITEUR ---
\noindent
\begin{minipage}[t]{0.5\textwidth}
    \textbf{Loïc Aron MBASSI EWOLO}\\
    Élève-Ingénieur Bac+5 -- Double diplôme ENSEA\\[3pt]
    \faMapMarker\ Cergy, France\\
    \faPhone\ (+33) 7 59 52 33 98\\
    \faEnvelope\ \href{mailto:loicaron.mbassiewolo@ensea.fr}{loicaron.mbassiewolo@ensea.fr}
\end{minipage}
\hfill
% --- DATE ---
\begin{minipage}[t]{0.4\textwidth}
    \raggedleft
    Cergy, le 5 février 2026
\end{minipage}

\vspace{1.2cm}

% --- DESTINATAIRE ---
\hfill
\begin{minipage}[t]{0.5\textwidth}
    \raggedleft
    \textbf{Groupe SII}\\
    Service Recrutement\\
    Toulouse (31), France
\end{minipage}

\vspace{1cm}

\noindent\textbf{Objet :} Candidature au stage Ingénieur Développement Logiciel Embarqué

\vspace{0.8cm}

Madame, Monsieur,

Actuellement élève-ingénieur en double diplôme à l'ENSEA, spécialisé en systèmes embarqués et automatique, je suis à la recherche d'un stage de fin d'études à partir d'avril 2026. Votre offre au sein du département Logiciel Embarqué, portant sur le développement de logiciels avioniques critiques, correspond parfaitement à mon projet professionnel et à mes compétences.

Le cycle en V et les exigences de certification du domaine aéronautique trouvent un écho direct dans mon parcours. Mon projet d'option à l'ENSEA sur la détection de défauts et le contrôle tolérant aux pannes pour drone m'a confronté aux enjeux de la sûreté de fonctionnement : modélisation en conditions nominales et dégradées, conception d'observateurs, validation par simulation et analyse de robustesse. Cette approche méthodique, de la spécification à la validation, est celle que vous attendez pour vos logiciels critiques.

Mon expérience en développement C/C++ embarqué s'est consolidée dans le cadre du Challenge CoHoMa (robotique autonome militaire), où je développe du code temps réel sur Nvidia Jetson avec des contraintes de fiabilité strictes. Par ailleurs, mon projet PROJET-TAXI (simulation multiprocessus en C sous Linux) m'a permis de maîtriser la gestion mémoire, la synchronisation inter-processus et la rédaction de spécifications fonctionnelles.

Je suis rigoureux, autonome et j'apprécie le travail en équipe. La promesse « Let's Tech Together » du Groupe SII résonne avec ma vision de l'ingénierie : construire ensemble des solutions innovantes dans un cadre exigeant. Je serais honoré de contribuer aux activités de conception, développement et validation de vos logiciels embarqués avioniques.

Je reste à votre disposition pour un entretien afin de discuter de ma candidature.

Je vous prie d'agréer, Madame, Monsieur, l'expression de mes salutations distinguées.

\vspace{0.8cm}

\textbf{Loïc Aron MBASSI EWOLO}\\
\textit{Élève-Ingénieur ENSEA / Polytechnique Yaoundé}

\end{document}
